\exposetitle{XXI}{Root data}
\label{XXI}
\begin{center}\textit{by M.~Demazure}\end{center}
{\renewcommand{\thefootnote}{(0)}\nde{Version of 13/10/2024}}%
This expos\'e collects, in the absence of a suitable reference,{\renewcommand{\thefootnote}{(1)}\nde{For the results on root systems (\S\S~1--5), one may consult [BLie], Chap.~VI.}} known results on root data (= ``abstract'' root systems), most of which will be used in what follows.

\textbf{Notations.} One denotes by $\QQ_+$ the set of positive (or zero) rational numbers; one has $\ZZ\cap\QQ_+=\NN$. Let $V$ be a $\QQ$-vector space; if $A$ (resp. $B$) is a subset of $\QQ$ (resp. $V$), one denotes by $A\cdot B$ the image of $A\otimes B$ by the morphism $\QQ\otimes V\to V$, in other words the set of linear combinations of elements of $B$ with coefficients in $A$. One puts $-B=\{-1\}B$. One denotes by $E-F$ the set of elements of $E$ which do not belong to $F$.

\sgasection{1}{Generalities}\label{XXI.sec.1}
\sgasection{1.1}{Definitions, first properties}
\begin{definition}{1.1.1}\label{XXI.1.1.1}
Let $M$ and $M^*$ be two free $\ZZ$-modules of finite type in duality. Put $V=M\otimes\QQ$, $V^*=M^*\otimes\QQ$; these are two $\QQ$-vector spaces in duality. Identify $M$ (resp. $M^*$) with a subset of $V$ (resp. $V^*$). The canonical bilinear form on $M^*\times M$ (resp. $V^*\times V$) is denoted $(\ ,\ )$.

Let $R$ be a finite subset of $M$. Give a map $\alpha\mto\alpha^*$ from $R$ to $M^*$; the set of $\alpha^*$ for $\alpha\in R$ is denoted $R^*$. With each $\alpha\in R$, associate the endomorphism $s_\alpha$ (resp. $s_\alpha^*$) of $M$ and $V$ (resp. $M^*$ and $V^*$) given by the formulas:
\begin{equation}\tag{1}\label{XXI.eq.1}
s_\alpha(x)=x-(\alpha^*,x)\alpha,\qquad\text{i.e. }s_\alpha=\mathrm{id}-\alpha^*\otimes\alpha;
\end{equation}
\begin{equation}\tag{$1^*$}\label{XXI.eq.1star}
s_\alpha^*(u)=u-(u,\alpha)\alpha^*,\qquad\text{i.e. }s_\alpha^*=\mathrm{id}-\alpha\otimes\alpha^*.
\end{equation}
One says that the pair $(R,R^*)$ (more precisely the pair $(R,R\to M^*)$) is a \emph{root datum} in $(M,M^*)$, or that $(M,M^*,R,R^*)$ is a root datum, if the following axioms hold:
\[
\begin{array}{ll}
\text{(DR I)}&\text{For each }\alpha\in R,\text{ one has }(\alpha^*,\alpha)=2.\\
\text{(DR II)}&\text{For each }\alpha\in R,\text{ one has }s_\alpha(R)\subset R,\quad s_\alpha^*(R^*)\subset R^*.
\end{array}
\]
One says that $R$ is the \emph{root system} of the root datum $\mathcal R=(M,M^*,R,R^*)$. The elements of $R$ (resp. $R^*$) are called the \emph{roots} (resp. \emph{coroots}) of the root datum.
\end{definition}

\begin{remark}{1.1.2}\label{XXI.1.1.2}
Axiom (DR I) is equivalent to any one of the following properties:
\[
\begin{array}{lll}
(2)\ s_\alpha s_\alpha=\mathrm{id},&&(2^*)\ s_\alpha^*s_\alpha^*=\mathrm{id},\\
(3)\ s_\alpha(\alpha)=-\alpha,&&(3^*)\ s_\alpha^*(\alpha^*)=-\alpha^*.
\end{array}
\]
\end{remark}
\begin{remark}{1.1.3}\label{XXI.1.1.3}
Axioms (DR I) and (DR II) imply
\[
R=-R,\qquad R^*=-R^*,\qquad0\notin R,\qquad0\notin R^*.
\]
\end{remark}
\begin{lemma}{1.1.4}\label{XXI.1.1.4}
The map $R\to R^*$ is a bijection. More generally, if $\alpha,\beta\in R$ and $(\alpha^*,x)=(\beta^*,x)$ for every $x\in R$, then $\alpha=\beta$.
\end{lemma}
Indeed, one then has $s_\beta(\alpha)=\alpha-2\beta$, $s_\alpha(\beta)=\beta-2\alpha$. One deduces at once
\[
\begin{gathered}
s_\beta s_\alpha(\alpha)=2\beta-\alpha=\alpha+2(\beta-\alpha),\\
s_\beta s_\alpha(\beta-\alpha)=s_\beta(\beta-\alpha)=\beta-\alpha,
\end{gathered}
\]
whence $(s_\beta s_\alpha)^n(\alpha)=\alpha+2n(\beta-\alpha)\in R$ by (DR II). Since $R$ is finite, one has $\beta-\alpha=0$.

\begin{corollary}{1.1.5}\label{XXI.1.1.5}
The inverse map $R^*\to R$ defines a root datum
\[
\mathcal R^*=(M^*,M,R^*,R)
\]
called the \emph{dual} of $\mathcal R$.{\renewcommand{\thefootnote}{(2)}\nde{One will note that $(\alpha^*)^*=\alpha$.}}
\end{corollary}

\begin{definition}{1.1.6}\label{XXI.1.1.6}
Denote by $\Gamma_0(R)$ the subgroup of $M$ generated by $R$. Denote by $\mathcal V(R)$ the vector subspace of $V$ generated by $R$, that is, $\Gamma_0(R)\otimes\QQ$. Applying these definitions to $\mathcal R^*$, one similarly constructs $\Gamma_0(R^*)$ and $\mathcal V(R^*)$.

One calls the \emph{reductive rank} of $\mathcal R$ the number
\[
\operatorname{rgred}(\mathcal R)=\operatorname{rank}(M)=\dim(V)=\dim(V^*)=\operatorname{rank}(M^*)=\operatorname{rgred}(\mathcal R^*).
\]
One calls the \emph{semisimple rank} of $\mathcal R$ the number
\[
\operatorname{rgss}(\mathcal R)=\operatorname{rank}(R)=\operatorname{rank}(\Gamma_0(R))=\dim(\mathcal V(R)).
\]
One therefore has $\operatorname{rgss}(\mathcal R)\leq\operatorname{rgred}(\mathcal R)$.

We shall see below that $\operatorname{rgss}(\mathcal R)=\operatorname{rgss}(\mathcal R^*)$, i.e. that $\mathcal V(R)$ and $\mathcal V(R^*)$ have the same dimension.
\end{definition}
\begin{definition}{1.1.7}\label{XXI.1.1.7}
One says that $\mathcal R$ is \emph{semisimple} (resp. \emph{trivial}) if $\operatorname{rgss}(\mathcal R)=\operatorname{rgred}(\mathcal R)$ (resp. $\operatorname{rgss}(\mathcal R)=0$). For $\mathcal R$ to be trivial, it is therefore necessary and sufficient that $R$ be empty. The trivial root datum of reductive rank zero is denoted $0=(\{0\},\{0\},\varnothing,\varnothing)$.
\end{definition}

\begin{definition}{1.1.8}\label{XXI.1.1.8}
Denote by $W(\mathcal R)$ the group of transformations of $M$ generated by the $s_\alpha$, $\alpha\in R$. It is called the \emph{Weyl group} of $\mathcal R$. Put
\[
W^*(\mathcal R)=W(\mathcal R^*).
\]
Then $W(\mathcal R)$ acts on $R$, $\Gamma_0(R)$, $\mathcal V(R)$, $M$ and $V$. If $w\in W(\mathcal R)$ and $x\in M$ (resp. $x\in V$), one has $wx-x\in\Gamma_0(R)$ (resp. $wx-x\in\mathcal V(R)$); this is immediate from formula (1). Similarly for $W^*(\mathcal R)$.
\end{definition}

\begin{lemma}{1.1.9}\label{XXI.1.1.9}
For all $\alpha\in R$, $x\in V$, $u\in V^*$, one has
\begin{equation}\tag{4}\label{XXI.eq.4}
(s_\alpha^*(u),s_\alpha(x))=(u,x).
\end{equation}
\end{lemma}
Indeed, taking the scalar product of (1) and ($1^*$), one finds that the left-hand side is equal to $(u,x)+(u,\alpha)(\alpha^*,x)((\alpha^*,\alpha)-2)=(u,x)$.

\begin{remark}{1.1.10}\label{XXI.1.1.10}
If one supposes $0\notin R$ and $0\notin R^*$, then 1.1.9 is equivalent to (DR I).
\end{remark}
\begin{corollary}{1.1.11}\label{XXI.1.1.11}
Let $h\mto h^\vee$ denote the isomorphism from $\GL(M)$ onto $\GL(M^*)$ associating with $h$ its contragredient. Then formula (4) is also written
\begin{equation}\tag{5}\label{XXI.eq.5}
(s_\alpha)^\vee=s_\alpha^*.
\end{equation}
\end{corollary}
\begin{corollary}{1.1.12}\label{XXI.1.1.12}
The preceding isomorphism induces an isomorphism from $W(\mathcal R)$ onto $W^*(\mathcal R)$.
\end{corollary}
\begin{scholium}{1.1.13}\label{XXI.1.1.13}
By the preceding result, we shall identify $W$ and $W^*$, and consider $W$ as a group of transformations of $R$, $R^*$, $M$, $M^*$, $\Gamma_0(R)$, $\Gamma_0(R^*)$, $V$, $V^*$, $\mathcal V(R)$, $\mathcal V(R^*)$. We shall write $s_\alpha$ for $s_\alpha^*$.
\end{scholium}

\sgasection{1.2}{The map $p$}
\begin{lemma}{1.2.1}\label{XXI.1.2.1}
Let $p:M\to M^*$ (resp. $V\to V^*$) be the linear map defined by
\begin{equation}\tag{6}\label{XXI.eq.6}
p(x)=\sum_{u\in R^*}(u,x)u.
\end{equation}
Put $\ell(x)=(p(x),x)$. One has the following properties:
\begin{equation}\tag{7}\label{XXI.eq.7}
\ell(x)\geq0,\qquad\ell(\alpha)>0\quad\text{for }\alpha\in R.
\end{equation}
\begin{equation}\tag{8}\label{XXI.eq.8}
\ell(wx)=\ell(x),\qquad\text{for }w\in W.
\end{equation}
\begin{equation}\tag{9}\label{XXI.eq.9}
(p(x),y)=(p(y),x),\qquad\text{for }x,y\in V.
\end{equation}
\begin{equation}\tag{10}\label{XXI.eq.10}
\ell(\alpha)\alpha^*=2p(\alpha),\qquad\text{for }\alpha\in R.
\end{equation}
\end{lemma}
The first three relations are immediate.{\renewcommand{\thefootnote}{(3)}\nde{(8) follows from (4), (2) and (DR II).}} Let us prove the last. By (1), for $u\in V^*$, one has:
\[
\begin{aligned}
(u,\alpha)^2\alpha^*&=(u,\alpha)u-(u,\alpha)s_\alpha(u)\\
&=(u,\alpha)u+(u,s_\alpha(\alpha))s_\alpha(u)\\
&=(u,\alpha)u+(s_\alpha(u),\alpha)s_\alpha(u)\qquad\text{(since }s_\alpha^2=\mathrm{id}\text{).}
\end{aligned}
\]
Since $s_\alpha$ is a permutation of $R^*$ (by (DR II)), one need only sum over $u\in R^*$ to conclude.

\begin{scholium}{1.2.2}\label{XXI.1.2.2}
Relation (10) says that $\alpha\mto\ell(\alpha)\alpha^*$ is the restriction to $R$ of a linear map from $M$ to $M^*$. In particular, one has
\[
(-\alpha)^*=-\alpha^*.
\]
\end{scholium}
\begin{corollary}{1.2.3}\label{XXI.1.2.3}
The map $p$ induces an isomorphism from $\mathcal V(R)$ onto $\mathcal V(R)^*$.
\end{corollary}
Indeed, $p$ sends $\mathcal V(R)$ onto $\mathcal V(R^*)$. One therefore has
\[
\dim(\mathcal V(R))\geq\dim(\mathcal V(R^*)).
\]
Applying this formula to the dual root datum, one deduces
\[
\dim(\mathcal V(R))=\dim(\mathcal V(R^*)),
\]
hence $p$, being surjective, is also bijective.

\begin{corollary}{1.2.4}\label{XXI.1.2.4}
One has $\dim(\mathcal V(R))=\dim(\mathcal V(R^*))$, hence
\[
\operatorname{rgss}(\mathcal R)=\operatorname{rgss}(\mathcal R^*).
\]
\end{corollary}
\begin{corollary}{1.2.5}\label{XXI.1.2.5}
The bilinear form $(u,x)$ is nondegenerate on $\mathcal V(R^*)\times\mathcal V(R)$, hence puts these $\QQ$-vector spaces in duality.
\end{corollary}
Indeed, if $(u,x)=0$ for every $u\in R^*$, then $p(x)=0$.
\begin{corollary}{1.2.6}\label{XXI.1.2.6}
The symmetric bilinear form $(p(x),y)$ is positive nondegenerate on $\mathcal V(R)$.
\end{corollary}
\begin{corollary}{1.2.7}\label{XXI.1.2.7}
$W$ acts faithfully on $R$ (and hence on the other sets of 1.1.13).
\end{corollary}
Indeed, let $u\in V^*$. Let $w\in W$; suppose $w(\alpha)=\alpha$ for every $\alpha\in R$, and let us prove $w(u)=u$. One has
\[
(w(u)-u,\alpha)=(w(u),\alpha)-(u,\alpha)=(u,w^{-1}(\alpha))-(u,\alpha)=0.
\]
But $w(u)-u\in\mathcal V(R^*)$. If it is orthogonal to all the roots, it is zero by 1.2.5.
\begin{corollary}{1.2.8}\label{XXI.1.2.8}
The group $W$ is finite.
\end{corollary}

\begin{proposition}{1.2.9}\label{XXI.1.2.9}
The actions of $W$ respect the correspondence between roots and coroots. In other words, for $\alpha\in R$ and $w\in W$, one has
\[
w(\alpha^*)=w(\alpha)^*.
\]
\end{proposition}
It suffices to check this for $w=s_\beta$, $\beta\in R$, i.e. to check the formula
\[
s_\beta(\alpha^*)=s_\beta(\alpha)^*.
\]
Now $\ell(s_\beta(\alpha))s_\beta(\alpha)^*/2$ equals:
\[
\begin{aligned}
p(s_\beta(\alpha))&=\sum_{u\in R^*}(u,s_\beta(\alpha))u=\sum_{u\in R^*}(s_\beta(u),\alpha)u\\
&=\sum_{u\in R^*}(u,\alpha)s_\beta(u)=s_\beta(p(\alpha));
\end{aligned}
\]
since $\ell(s_\beta(\alpha))=\ell(\alpha)$, one obtains $s_\beta(\alpha)^*=s_\beta(2p(\alpha)/\ell(\alpha))=s_\beta(\alpha^*)$.

\begin{corollary}{1.2.10}\label{XXI.1.2.10}
If $\alpha\in R$ and $w\in W$, one has
\[
ws_\alpha w^{-1}=s_{w(\alpha)}.
\]
\end{corollary}
Indeed, $ws_\alpha w^{-1}(x)=x-(\alpha^*,w^{-1}(x))w(\alpha)=x-(w(\alpha^*),x)w(\alpha)$, and this last term equals, by 1.2.9:
\[
x-(w(\alpha)^*,x)w(\alpha)=s_{w(\alpha)}(x).
\]
\begin{corollary}{1.2.11}\label{XXI.1.2.11}
Let $R'\subset R$. For $(M,M^*,R',R'^*)$ to be a root datum, it is necessary and sufficient that $\alpha,\beta\in R'$ imply $s_\alpha(\beta)\in R'$.
\end{corollary}

\sgasection{2}{Relations between two roots}\label{XXI.sec.2}
\sgasection{2.1}{Proportional roots}
\begin{proposition}{2.1.1}\label{XXI.2.1.1}
Let $\alpha$ and $\beta$ be two roots. The following conditions are equivalent:
\begin{enumeratei}
\item There exists $k\in\QQ$ such that $\alpha=k\beta$.
\item $s_\alpha=s_\beta$.
\end{enumeratei}
Moreover, under these conditions, one has $\alpha^*=k^{-1}\beta^*$, and $k$ is equal to one of the numbers $1$, $-1$, $2$, $-2$, $1/2$, $-1/2$.
\end{proposition}
Suppose first (i). First, one has
\[
\alpha^*=\ell(\alpha)^{-1}2p(\alpha)=k^{-2}\ell(\beta)^{-1}2kp(\beta)=k^{-1}\beta^*.
\]
This implies at once $s_\alpha=s_\beta$. Conversely, if one has $s_\alpha=s_\beta$, then
\[
\alpha=s_\alpha(-\alpha)=s_\beta(-\alpha)=(\beta^*,\alpha)\beta-\alpha,
\]
whence $2\alpha=(\beta^*,\alpha)\beta$, with $(\beta^*,\alpha)\in\ZZ$, hence (ii) implies (i). Finally, if $\alpha=k\beta$, then $\alpha^*=k^{-1}\beta^*$, whence
\[
(\alpha^*,\beta)=2k^{-1},\qquad(\beta^*,\alpha)=2k,
\]
hence $2k$ and $2k^{-1}$ are integers, and the proof is finished.

\begin{enoncerm}{Application}{2.1.2}\label{XXI.2.1.2}
Root data $\mathcal R$ such that $\operatorname{rgss}(\mathcal R)=1$ are of one of the following two types:
\begin{enumeratei}
\item Type $A_1$: there exists $\alpha\in M$ such that the roots are $\alpha$ and $-\alpha$. The coroots are then $\alpha^*$ and $-\alpha^*$.
\item Type $A'_1$: there exists $\alpha\in M$ such that the roots are $\alpha$, $-\alpha$, $2\alpha$, $-2\alpha$. The coroots are then $\alpha^*$, $-\alpha^*$, $\alpha^*/2$, $-\alpha^*/2$.
\end{enumeratei}
\end{enoncerm}

\begin{definition}{2.1.3}\label{XXI.2.1.3}
One says that $\alpha\in R$ is \emph{indivisible} if $\alpha/2\notin R$. One says that $\mathcal R$ is \emph{reduced} if every root is indivisible.
\end{definition}
For $\mathcal R$ to be reduced, it is necessary and sufficient that $\mathcal R^*$ be so. If $\alpha$ is indivisible and $w\in W$, then $w(\alpha)$ is indivisible.
\begin{definition}{2.1.4}\label{XXI.2.1.4}
Let $\alpha\in R$. If $\alpha$ is indivisible, put $\operatorname{ind}(\alpha)=\alpha$. Otherwise, put $\operatorname{ind}(\alpha)=\alpha/2$.
\end{definition}
\begin{corollary}{2.1.5}\label{XXI.2.1.5}
If $\alpha\in R$ is indivisible and $k\alpha\in R$, where $k\in\QQ$, then $k\in\ZZ$.
\end{corollary}
\begin{proposition}{2.1.6}\label{XXI.2.1.6}
Let $\mathcal R$ be a root datum. Then
\[
\operatorname{ind}(\mathcal R)=(M,M^*,\operatorname{ind}(R),\operatorname{ind}(R)^*)
\]
is a reduced root datum, and one has
\[
W(\operatorname{ind}(\mathcal R))=W(\mathcal R).
\]
\end{proposition}
Indeed, $\operatorname{ind}(\mathcal R)$ is a root datum by 1.2.11, since the Weyl group permutes the indivisible roots. The second assertion follows from 2.1.1.
\begin{remark}{2.1.7}\label{XXI.2.1.7}
If $\mathcal R$ is not reduced, one has $\operatorname{ind}(R)^*\ne\operatorname{ind}(R^*)$, and hence $\operatorname{ind}(\mathcal R^*)\ne\operatorname{ind}(\mathcal R)^*$.
\end{remark}

\sgasection{2.2}{Orthogonal roots}
\begin{lemma}{2.2.1}\label{XXI.2.2.1}
Let $\alpha$ and $\beta$ be two roots. One has
\begin{equation}\tag{11}\label{XXI.eq.11}
\ell(\alpha)(\alpha^*,\beta)=\ell(\beta)(\beta^*,\alpha).
\end{equation}
\end{lemma}
This follows at once from 1.2.1, formulas (9) and (10).
\begin{corollary}{2.2.2}\label{XXI.2.2.2}
Let $\alpha,\beta\in R$. The following conditions are equivalent:
\begin{enumeratei}
\item $(\alpha^*,\beta)=0$,\qquad(i bis) $(\beta^*,\alpha)=0$,
\item $(p(\alpha),\beta)=0$,
\item $s_\alpha(\beta)=\beta$,\qquad(iii bis) $s_\beta(\alpha)=\alpha$,
\item $s_\alpha(\beta^*)=\beta^*$,\qquad(iv bis) $s_\beta(\alpha^*)=\alpha^*$,
\item $s_\alpha\ne s_\beta$, and $s_\alpha$ and $s_\beta$ commute.
\end{enumeratei}
\end{corollary}
All the equivalences are immediate except those involving (v). Let us show that (i) (and (i bis)) imply (v). One has
\[
s_\alpha s_\beta(x)=x-(\beta^*,x)\beta-(\alpha^*,x)\alpha+(\beta^*,x)(\alpha^*,\beta)\alpha.
\]
If $(\alpha^*,\beta)=0$, then $(\beta^*,\alpha)=0$, and $s_\alpha s_\beta(x)=s_\beta s_\alpha(x)$.

Conversely, suppose (v). One then has
\[
s_\alpha=s_\beta s_\alpha s_\beta=s_{s_\beta(\alpha)}\qquad\text{(by 1.2.10).}
\]
By 2.1, there therefore exists $k\in\QQ^*$ such that $\alpha=ks_\beta(\alpha)=k\alpha-k(\beta^*,\alpha)\beta$. Since $s_\alpha\ne s_\beta$, $\beta$ and $\alpha$ are not proportional by 2.1.1, hence $(\beta^*,\alpha)=0$.
\begin{definition}{2.2.3}\label{XXI.2.2.3}
Two roots satisfying the equivalent conditions of 2.2.2 are called \emph{orthogonal}.
\end{definition}
